%! AP Statistics — Unit 9, Lesson 9.1 — Exit Ticket (Answer Key)
\documentclass[10pt]{article}
\usepackage{apstats-article}
\usepackage{apstats-key}

\begin{document}

\pageheader{Unit 9, Lesson 9.1}{Exit Ticket --- Answer Key}
\namedateperiod

\vspace{0.1in}
A researcher collects data from 12 stores on the number of ads run per month ($x$) and
monthly revenue (\$thousands, $y$). The least-squares regression line is
$\hat{y} = 45.2 + 0.8x$.

\begin{enumerate}

  \item Interpret the slope $b = 0.8$ in context.

        \ans{For each additional ad run per month, the predicted monthly revenue
        increases by \$0.8 thousand (or \$800).}

  \item Another researcher says, ``Since $b \ne 0$, we can conclude that advertising
        increases revenue.'' What is wrong with this reasoning?

        \ans{The slope $b = 0.8$ is a sample statistic. Even if $\beta = 0$ (advertising
        has no real effect on revenue), random variation in the data could produce a
        non-zero $b$. We cannot conclude $\beta \ne 0$ just because $b \ne 0$.}

  \item What question does Unit 9 ask about this slope, and why can't we answer it just
        by looking at $b$?

        \ans{Unit 9 asks whether the true population slope $\beta$ is different from
        $0$ --- i.e., whether there is a real linear relationship between advertising and
        revenue in the population. We cannot answer this just from $b$ because sampling
        variability always produces non-zero slopes, even when $\beta = 0$.}

\end{enumerate}

\end{document}
